\chapter{Proposed Constricted-Channel Study}\label{ch:constriction}

\section{Purpose and status of the extension}
A constriction introduces a varying gap, a changing cross-sectional mean
velocity at fixed flux, and stronger sensitivity to wall treatment.
It provides a natural next problem after the constant-width baseline and
the Gaussian coupling study. It also introduces new requirements that
cannot be met by merely changing the appearance of a straight-channel
figure.

This chapter defines the geometry, governing problem, and evaluation design
for that extension. The geometry drawings are mathematical illustrations.
No constricted-channel velocity, pressure, particle trajectory, or deposition
rate has been computed for the accepted result set. The proposed parameter
values below are study-design choices rather than measured or completed
results.

The obstruction is static and represented by the walls. A progressively
growing particle deposit, moving obstacle, or porous plug would require
additional evolution laws and is outside this design. I use the term
\emph{constriction} for the prescribed narrowing to keep this geometric
effect separate from dynamic clogging.

\section{Parameterized geometry}
Let the nominal channel height be $H$ and let the walls be
$y=\pm h(x)$. A compact, smooth narrowing centered at $x_c$ is defined by
\begin{equation}
 h(x)=\frac{H}{2}[1-\delta f(\xi)],\qquad
 \xi=\frac{x-x_c}{\ell},\qquad 0\leq\delta<1,
 \label{eq:constriction_height}
\end{equation}
where
\begin{equation}
 f(\xi)=
 \begin{cases}
 \left[\dfrac{1+\cos(\pi\xi)}{2}\right]^2,&|\xi|<1,\\[2mm]
 0,&|\xi|\geq1.
 \end{cases}
 \label{eq:constriction_profile}
\end{equation}
The profile and its first three derivatives match the constant-height
segments at the joins. Its minimum gap is $H_t=(1-\delta)H$. The initial
design uses $L_x/H=8$, $x_c/H=4$, and $\ell/H=1.2$, with nominal severities
$\delta\in\{0,0.2,0.4\}$. These values define a controlled sequence from
the unblocked channel to moderate narrowing. The actual reference study
must check that upstream and downstream lengths are sufficient.

\Cref{fig:constriction_geometry} displays the geometry and height profiles.
Only geometric quantities are plotted. In particular, the height curves
must not be interpreted as velocity or pressure results.

\ResultFigure{13_constriction_geometry}{Proposed constriction. Panel (a)
defines the throat, characteristic length, and reference sections. Panel
(b) shows the mathematical height profiles for the proposed severity
values. The illustrated channel has $L_x/H=8$ and $\ell/H=1.2$; no fluid
or particle solution is represented.}{fig:constriction_geometry}

\section{Carrier-flow boundary conditions}
For an open-channel extension, the governing equations retain the inertial
pressure scale and the same Newtonian constitutive law. I prescribe a
parabolic inlet velocity with a specified flux $Q$ and no slip on both
walls. A downstream traction condition provides an outlet reference rather
than imposing a second velocity profile. The finite-element reference and
the PINN must implement exactly the same outlet condition and pressure
convention.

This is a different boundary-value problem from both earlier stages.
The curved baseline has prescribed velocities at both ends, whereas the
straight Gaussian channel is periodic and body-force driven. Comparisons
across these stages must therefore identify which changes follow from
geometry and which follow from the drive or boundary conditions. A useful
control is a straight open channel with the same inlet and outlet treatment
as the constricted case.

For unsteady particle introduction, the initial carrier field should be an
independently converged unladen solution in the constricted domain.
Imposing the original straight-channel profile throughout the narrowed
region would violate the new wall geometry. An alternative start from rest
requires a compatible smooth inlet ramp and its own transient verification.
The primary design uses the steady unladen field to avoid mixing inlet
startup with particle response.

\section{Geometry-aware representation}
The fixed-width centerline mapping of the baseline cannot represent this
variable-height domain without modification. A convenient mapping is
\begin{equation}
 (x,\eta)\longmapsto(x,y)=(x,h(x)\eta),\qquad -1\leq\eta\leq1.
 \label{eq:constriction_map}
\end{equation}
Its area Jacobian is $h(x)$. If the differential residual is evaluated in
mapped coordinates, the chain rule must include this geometry dependence.
For example,
\begin{equation}
 \left.\partial_x\right|_y
 =\left.\partial_x\right|_\eta
 -\eta\frac{h'(x)}{h(x)}\partial_\eta,\qquad
 \partial_y=\frac{1}{h(x)}\partial_\eta.
 \label{eq:mapped_derivatives}
\end{equation}
Evaluating automatic derivatives directly with respect to physical
coordinates is another option. Either route requires a manufactured
derivative test before training a new flow problem.

Periodic Fourier features in the streamwise direction would impose the
wrong structure for the proposed open domain. They must be replaced by
nonperiodic coordinate features. A wall-distance factor can enforce the
stationary wall velocity, but it must be consistent with the nonzero inlet
condition and with the mapped derivatives \citep{sukumar2022}. % Citation-plan C27
 A divergence-free streamfunction
ansatz is a possible alternative, provided that its boundary and pressure
conditions are verified. These are implementation choices to evaluate,
rather than capabilities already established by the straight-channel model.

\section{Gaussian support near the narrowing}
The force-transfer rule must use quadrature over the fluid domain:
\begin{equation}
 Z_i=\sum_{a\in\Omega_\delta}
 G(\vect{x}_a-\vect{q}_i)\Delta A_a,\qquad
 W_{ai}=\frac{G(\vect{x}_a-\vect{q}_i)}{Z_i}.
 \label{eq:constriction_normalization}
\end{equation}
For mapped quadrature, $\Delta A_a$ includes the Jacobian from
\cref{eq:constriction_map}. The same points, weights, and normalizer must
be used for interpolation and reaction-force deposition. Under that
condition, the derivations of \cref{eq:force_identity,eq:work_identity}
remain valid on the new quadrature.

\Cref{fig:boundary_kernel} illustrates why normalization must follow the
actual domain. Some Gaussian support intersects the wall region near the
throat. Removing those points without updating $Z_i$ would lose force.
Renormalizing restores the integrated force, but changes the shape of the
sampling operator, an issue related to coarse-graining choices
\citep{sun2015}. Physical wall corrections \citep{goldman1967,brenner1961}
and undisturbed-velocity recovery \citep{horwitz2016,ireland2017} remain
separate modeling requirements. % Citation-plan C28

\ResultFigure[.95]{14_boundary_kernel}{Illustration of a Gaussian centered
near the upper throat wall. Shading is a prescribed kernel shape restricted
to the proposed fluid domain, not a computed concentration or force field.
The normalizer must use the same fluid-domain quadrature as both transfer
operations.}{fig:boundary_kernel}

Three ratios require explicit monitoring:
\begin{equation}
 \frac{d}{H_t},\qquad \frac{\sigma}{H_t},\qquad
 \frac{\sigma}{\Delta_{\mathrm{quad}}},
 \label{eq:constriction_ratios}
\end{equation}
where $\Delta_{\mathrm{quad}}$ denotes a local quadrature-spacing measure.
The first concerns physical confinement; the second concerns regularization
relative to the gap; the third concerns the numerical resolution of that
regularization. They address different questions and should be varied
separately, with coarse-graining and transfer-resolution studies providing
related context \citep{sun2015,capecelatro2013}. % Citation-plan C29
 If physical wall effects dominate, the present unbounded
Stokes closure may need replacement before more accurate training is useful.

\section{Particle and parcel wall conditions}
An individual sphere has an admissible center region obtained by offsetting
the wall by its radius. For a sloped wall, the offset is based on signed
normal distance to the boundary, not only the vertical gap. A contact event
must therefore use the actual geometry. The baseline's irreversible
sticking rule can be retained as a controlled modeling choice, but it must
not be called a prediction of deposition physics.

A parcel introduces an additional issue: its center represents many
particles. Stopping the entire represented population at the first center
contact may create a strong and artificial capture rule. The proposed
first study therefore focuses on carrier fields and freely transported
dilute parcels whose centers remain inside the admissible region. A
deposition extension should specify whether it resolves individual
particles, splits parcels, or uses a separately justified population-level
wall model. No capture efficiency is assigned before making that choice.

\section{Controlled evaluation sequence}
The extension uses staged verification consistent with the distinction
between solving equations and establishing physical validity
\citep{oberkampf2002}. The following sequence is the study design adopted
here. % Citation-plan C30
\begin{enumerate}
\item Recover the straight, unladen open-channel solution at $\delta=0$
with the new inlet and outlet conditions. Verify pressure gauge, wall
velocity, and section flux before adding particles.
\item Solve the unladen constricted channel with an independent mesh-based
method. Refine the throat and the upstream and downstream transition
regions; monitor pressure drop and the complete velocity field.
\item Introduce individual one-way particles with identical initial
conditions in the neural and reference fields. Compare position,
common-time coverage, and endpoint class separately.
\item Add a single Gaussian parcel at low loading and verify force/work
transfer on the irregular-domain quadrature. Compare the coupled fluid
response with an independently discretized regularized model.
\item Vary severity and loading only after the earlier controls pass.
Keep physical diameter, total represented mass, and kernel width explicit.
\end{enumerate}

The initial proposed Reynolds numbers are $1$, $5$, and $20$ at nominal
height. These are design points rather than an observed regime map.
Particle response times should be selected to retain an acceptably small
slip Reynolds number at the throat. Mass-loading changes must be made
through represented population while holding diameter fixed when the
purpose is to isolate feedback. A separate diameter sweep changes both
response time and confinement, and should be interpreted accordingly.

\section{Observables and acceptance criteria}
The primary observables are the section flux, throat velocity profile,
pressure drop between fixed reference sections, and field error against
the matched reference. For particles, residence time and endpoint class
are added only after trajectory agreement is demonstrated. The reference
sections must remain identical across severity values so that the pressure
drop represents the same geometric interval.

At fixed flux, incompressibility gives the exact sectional mean relation
\begin{equation}
 \overline{u}_x(x)=\frac{Q}{2h(x)},\qquad
 \frac{\overline{u}_{x,t}}{\overline{u}_{x,0}}
 =\frac{1}{1-\delta}.
 \label{eq:mean_acceleration}
\end{equation}
This is a geometric conservation implication, not a computed centerline
velocity or a prediction of recirculation. A lubrication estimate,
derived in \cref{app:derivations}, gives a useful additional check only
when its small-slope and low-inertia assumptions apply.

\begin{table}[htbp]\centering\small
\caption{Proposed acceptance protocol for the constricted study. Limits
must be fixed before interpreting the associated parameter sweep; these
checks have not been executed.}\label{tab:constriction_protocol}
\begin{tabularx}{\textwidth}{@{}p{.26\textwidth}YY@{}}\toprule
Quantity & Proposed check & Interpretation\\\midrule
Unladen field & $E_u,E_p<5\%$ against a matched reference & Carrier-field agreement\\
Reference refinement & Changes in $\Delta p$ and field norms below 2\% & Numerical sensitivity, not exact uncertainty\\
Section flux & Maximum relative variation below 2\% & Integrated mass balance\\
Transfer & Force/work discrepancies consistent with roundoff & Algebraic operator consistency\\
Particle paths & RMS position below $0.05W$; common interval above 95\%; all endpoint classes matched & Transport accuracy for stated initial conditions\\
Physical regime & Report maximum $\Rey_p$, $d/H_t$, local $\alpha$ and wall distance & Closure applicability; no universal threshold inferred\\
\bottomrule\end{tabularx}
\end{table}

The error limits in \cref{tab:constriction_protocol} are proposed numerical
screens. In addition, residuals should be reported by component on a
separate set of throat, near-wall, and bulk points. A global average must
not conceal a concentration of error near the narrowing. The Gaussian
width study should remain separate from spatial and temporal refinement,
and multiple seeds should be included before drawing a quantitative trend.

\section{Expected mechanisms and testable statements}
Conservation at fixed flux predicts a larger section-mean speed as the
gap narrows. The accompanying changes in particle transit depend on
response time, local flow, and the retained force law. I treat changes in
residence time, wall approach, and downstream recovery as hypotheses for
the controlled study. Recirculation, enhanced deposition, or a particular
pressure-loss factor cannot be inferred from the geometry alone.

At fixed imposed pressure difference, flux becomes an outcome instead.
That experiment addresses a different question and must not be pooled with
the fixed-flux design. Keeping these operating conditions separate is
essential for a meaningful comparison of unblocked and constricted cases.

\section{Chapter conclusions}
The constricted-channel stage is defined at the study-design level through
its geometry, boundary-condition choices, transfer consistency, and
evaluation requirements. The outlet traction and its reference value must
be fixed before executing the matched solvers.
It extends the study in a controlled direction while preserving the
distinction between mathematical consequences, proposed experiments, and
completed numerical findings. Its next milestone is a verified unladen
reference problem, followed by one-way transport and then two-way Gaussian
feedback.
